%% how-to-use --- how the book is meant to be read, and what its boxes mean.
%%
%% Twin: frontmatter/pl/how-to-use.tex. parity.py compares them token by
%% token, like every chapter.

\chapter*{\lblHowToUse}
\phantomsection
\addcontentsline{toc}{chapter}{\lblHowToUse}
\markboth{\lblHowToUse}{\lblHowToUse}

\noindent
This book assumes you can add, multiply, work out a percentage and read a
graph. It does not assume you remember any algebra, and it never uses a
symbol it has not introduced first. Every formula is said in words as well,
and where a formula would only get in the way, it is left out. If you
stopped doing mathematics at school and have always wanted to know what
people mean when they say \enquote{chaos}, you are the reader it was written
for.

What it does ask of you is to take part. The book is built as a sequence of
small steps, and at the end of many of them you are asked a question.

\section*{The loop}

Every chapter is read the same way:

\begin{enumerate}[leftmargin=1.6em, itemsep=3pt]
  \item \textbf{Read a short stretch.} Each section is a few paragraphs.
  \item \textbf{Stop and think.} A violet box asks you a question. Answer it
        \dash{} on paper, aloud, or as a guess \dash{} before you read on.
        The answer is directly underneath, so cover it with your hand.
  \item \textbf{Run the code}, if you have a computer beside you. Every
        listing is a small file, and running it takes a single command.
  \item \textbf{Do the lab}, if you want to make the idea your own: a lab is
        a half-written program and a test that tells you when it works.
  \item \textbf{Take the checkpoint} at the end of the chapter, and check
        your answers in Appendix~\ref{app:answers}.
\end{enumerate}

\section*{Why the book keeps asking you questions}

Being asked a question before you are told the answer is uncomfortable.
It is meant to be. An answer you have reached for, even a wrong one, stays
with you far longer than one you have only read, and a wrong guess that is
then corrected is the fastest way to find out which of your intuitions are
going to mislead you. This subject is full of those. Most people's first
guess about how fast a small error grows, or about what a simple rule can do,
is wrong, and the book is arranged so that you make that guess before you are
shown why.

So when a box asks you to answer before you read on, please do. Nobody is
marking it.

\section*{The laboratory}

Everything in the book that a computer can check has been checked by one,
and you can check it again. The programs are in the book's repository, under
\texttt{code/}. They are written in Python and need one tool, \pkg{uv},
which installs the right version of Python and every library by itself.
Once you have it, from the repository's \texttt{code/} directory:

\begin{shellcmd}
uv sync
uv run python ch01/falling_stone.py
uv run pytest -k l01_01
\end{shellcmd}

The first line sets up the laboratory, once. The second runs a listing from
Chapter~\ref{ch:models}. The third runs the test of that chapter's first
lab. You do not need to know Python to begin: the first chapters explain
every line they print, and the programs stay short throughout.

\section*{The boxes}

Besides the questions, the text is broken up by boxes, each with one job:

\begin{description}[leftmargin=1.2em, itemsep=2pt]
  \item[\lblTrap.] A mistake almost everybody makes, set out after you have
    had the chance to make it yourself.
  \item[\lblWorld.] A real place the idea turns up: a field, a device, a
    practice.
  \item[\lblRigour.] What the book states without proving it, and where the
    proof lives.
  \item[\lblNote, \lblWarning.] An aside, and a caution.
  \item[\lblProject.] A piece of the book's own small library of chaos, which
    the listings share.
  \item[\lblLab.] Something to do at the keyboard, with a test that tells
    you when you have done it.
\end{description}

\section*{Numbers you can trust}

Any number in this book that you could not work out in your head was
produced by a program, and the book prints the program's output rather than
a number somebody typed. The build checks every one of them again on every
change. That matters more than usual here, for a reason the book spends a
chapter on: in a chaotic system a computation can go wrong in the last digit
and be completely wrong a few dozen steps later. Chapter~\ref{ch:computers}
explains how the book protects itself against that, and why some of its
numbers are deliberately rounded.

\section*{Reading without a computer}

You can. Every listing's output is printed beside it, every figure is on the
page, and the questions work on paper. The labs are the only part you will
miss, and they are there for readers who want to go further, not to gate
the argument.

\section*{Two languages, one book}

The book is published in English and in Polish from one source. The two
editions contain the same chapters, sections, figures, numbers and questions
in the same order, and the build compares them to make sure they do.
